\section{本讲主线：数据不是被“收集”出来的，而是被“加工”出来的}

官方标题是 \textbf{Lecture 14: Data II}。上一讲讨论了数据从哪里来：live services、Common Crawl、GitHub、论文、书籍、许可数据，以及版权、robots.txt、Terms of Service 等边界。本讲把镜头推进到管线内部：即使已经获得 raw data，它离训练样本仍然很远。真实系统需要依次解决格式转换、质量筛选、重复消除、分布混合和合成数据验证。核心问题因此从“有没有数据”转成“这条加工链最终制造了怎样的训练分布”。

把这条管线写成一个最小抽象，可以得到：

$$
\mathcal{D}_{\mathrm{train}}
=\operatorname{Mix}\!\left(
\operatorname{Dedup}\!\left(
\operatorname{Filter}\!\left(
\operatorname{Transform}(\mathcal{D}_{\mathrm{raw}})
\right)\right)\right).
$$

其中，$\mathcal{D}_{\mathrm{raw}}$ 表示网页、PDF、代码仓库等原始对象；$\operatorname{Transform}$ 把异构格式线性化为模型可消费的文本或序列；$\operatorname{Filter}$ 依据语言、质量、安全等目标打分；$\operatorname{Dedup}$ 消除 exact 或 near duplicate；$\operatorname{Mix}$ 决定各 domain 的采样权重。

\begin{importantbox}{本讲最重要的工程判断}
数据处理不是一串互相独立的清洗脚本。每一步都会改变后续步骤看到的分布：转换器丢掉表格结构，过滤器会放大 target data 的偏好，去重会改变高频模式，混合比例又会改变梯度来自哪些 domain。设计数据管线时必须把它看成一个联动系统。
\end{importantbox}

\teachervoice{课堂总结明确提醒，真实数据工作往往很 ``grungy''、高度 domain-specific，并且必须回到 concrete examples 才能做出高质量数据集；自动指标只能帮助定位，不能替代样本检查。}

\subsection{一张数据账本}

为了避免后面被工具名和论文名淹没，本节先建立一张统一账本。每一个阶段都要同时记录输入对象、优化目标、可调参数和失败模式；这样模型指标变化时，团队才能沿 provenance 反向定位，而不是把所有变化都归因于“数据质量”。

\begin{table}[H]
\centering
\begin{tabular}{L{0.14\textwidth}L{0.25\textwidth}L{0.25\textwidth}L{0.22\textwidth}}
\toprule
阶段 & 输入 & 关键决策 & 典型失败 \\
\midrule
Transform & HTML、PDF、repo、WARC & 保留哪些结构，如何 linearize & 导航污染、OCR 错误、代码文件顺序丢失 \\
Filter & 大量候选 document & target、score、threshold、sampling & domain bias、误杀长尾、classifier shortcut \\
Dedup & 文档或片段集合 & exact/fuzzy 粒度、相似度阈值 & 泄漏未消除、模板文本过度删除 \\
Mix & 多个 domain corpus & token budget 与 mixture weights & benchmark 提升但真实分布退化 \\
Synthetic & prompt、teacher、verifier、environment & 生成与验证闭环 & 自我强化错误、格式投机、污染评测 \\
\bottomrule
\end{tabular}
\caption{本讲的数据处理账本：每一步都同时有算法问题和治理问题。}
\end{table}

\begin{knowledgebox}{术语表：先把五个动作分清楚}
\begin{tabular}{L{0.18\textwidth}L{0.72\textwidth}}
\toprule
术语 & 在本讲中的精确定义 \\
\midrule
Transformation & 把网页、PDF、repository 等对象转换成可训练序列，同时记录结构损失和转换版本。 \\
Filtering & 用明确的 target、score 与 threshold 决定哪些样本进入候选训练集。 \\
Deduplication & 在指定粒度与相似度定义下，降低同一内容被重复计权或泄漏到评测的概率。 \\
Data mixing & 在固定 token budget 下，为不同 domain 分配采样权重，从而改变梯度来源。 \\
Synthetic data & 由 generator 构造任务或答案，并由 verifier 与 environment 提供可审计反馈的数据。 \\
\bottomrule
\end{tabular}
\end{knowledgebox}

\subsection{本章小结}

这节课的主线不是“有哪些数据集”，而是怎样把 raw objects 变成一个可解释的训练分布。后续所有细节都可以回到四个问题：输入是什么、目标是什么、算法如何近似、错误会怎样传到模型里。

